
\item 
\textbf{HydroSense Tech }    \hfill     
\hfill  \textit{Sep 2017 -- Present (Part Time, remote); 
\item  Co-founder / Machine Learning Engineer\hfill May 2025 -- Sep 2025 (Full Time, Singapore)}
\vspace*{-3pt}
\begin{itemize}[leftmargin=0.2in, topsep=1pt]
  \item Co-founded an IoT rainfall sensing startup and led ML for field-deployed devices, helping scale the product line from pilot deployments to \textbf{multi-million-dollar annual revenue} across municipal and industrial customers.
  \item Built an end-to-end deployment stack: C++ firmware on STM32-based controllers plus Python services on an edge gateway, with secure OTA updates over RS-485 (local bus) and Bluetooth/Wi-Fi (field apps), enabling \textbf{fleet-wide remote monitoring and configuration} for 1000+ devices.
  \item Designed and shipped an \textbf{on-device 1D-CNN} for temperature-drift compensation using PyTorch training + INT8 quantization, reducing weight-estimation MAE by \textbf{15\% in production} (20k+ installed sensors)
  \item \textbf{Patents:} Co-inventor on 3 patents, covering rainfall sensing hardware and ML-based drift compensation / EMI mitigation algorithms.
\end{itemize}


\vspace*{-1pt}
\item \textbf{Seno Medical}\\
\textit{Machine Learning Engineer} \hfill Buffalo, NY \quad \textit{May 2024 -- Aug 2024}
\vspace*{-3pt}
\begin{itemize}[leftmargin=0.2in, topsep=1pt]
  \item Collaborated with domain experts to build a \textbf{2k+ QA} dataset from  technical documents; fine-tuned a \textbf{LLaMA-based model (LoRA)} and benchmarked retrieval baselines (\textbf{BM25, MiniLM}) as part of a retrieval-augmented QA system.
  \item Designed an end-to-end \textbf{evaluation harness} with a label schema (factuality, coverage, style), automatic metrics, and Python error-analysis scripts; reduced hallucination rate from \textbf{30\% to 15\%} and increased \textbf{Recall@5 to 90\%}.
  \item Implemented a modular \textbf{RAG backend}: chunking and ingestion jobs, \textbf{FAISS}-based vector store, retriever / re-ranker layer, and LLM orchestration, served via \textbf{FastAPI} and containerized with \textbf{Docker} for on-prem deployment.
\item Set up Git-based workflows and basic \textbf{observability} (structured logs, latency / error-rate dashboards) and shipped reproducible Docker images plus deployment runbooks for downstream teams.


\end{itemize}





\vspace*{-1pt}
\item \textbf{Roswell Park Comprehensive Cancer Center}\\
\textit{Research Scientist} \hfill Buffalo, NY \quad \textit{May 2023 -- Aug 2023}
\vspace*{-3pt}
\begin{itemize}[leftmargin=0.2in, topsep=1pt]
   \item Designed a \textbf{multi-modal data pipeline} that aligns imaging with structured text reports, including de-identification, schema normalization, and automated QC, to produce clean image–text pairs for generative modeling.
  \item Built \textbf{dataset generation jobs} (filtering, sampling, augmentation, metadata tagging) to create balanced image–report datasets for conditional generation and downstream training.
  \item Implemented a \textbf{domain-specific latent diffusion model} (Stable Diffusion–style, DDIM/DPMS sampling), trained on \textbf{AWS}, for conditional image–to–text report generation and synthetic data augmentation.
  \item Automated \textbf{GPU training and experiment tracking} with standardized configs, logging, and comparison scripts, enabling fast iteration over model variants and dataset settings on the AWS stack. 
\end{itemize} 